\documentclass[11pt,a4paper]{amsart}
\usepackage{amsmath,amssymb,amsthm}
\usepackage{hyperref}
\usepackage{booktabs}
\usepackage{enumitem}

\newtheorem{theorem}{Theorem}
\newtheorem{lemma}[theorem]{Lemma}
\newtheorem{proposition}[theorem]{Proposition}
\newtheorem{corollary}[theorem]{Corollary}
\theoremstyle{definition}
\newtheorem{definition}[theorem]{Definition}
\newtheorem{remark}[theorem]{Remark}

\DeclareMathOperator{\SL}{SL}
\DeclareMathOperator{\PGL}{PGL}
\DeclareMathOperator{\tr}{tr}
\DeclareMathOperator{\diam}{diam}

\title[Zaremba's Conjecture]{Zaremba's Conjecture for $A=5$: \\
A Computer-Assisted Proof via Spectral Gaps and the Frolenkov--Kan Sieve}

\author{Cahlen Humphreys}
\address{Enfuse.io, Irvine, CA}
\email{cahlen@gmail.com}

\date{\today}

\begin{document}

\begin{abstract}
We prove Zaremba's Conjecture (1972) for the bound $A=5$: for every integer $d \geq 1$, there exists $a$ with $\gcd(a,d)=1$ such that all partial quotients of $a/d$ are at most~$5$. The proof is effective for all $d \leq 10^{1500}$ and extends to all $d$ via the non-effective property~$(\tau)$ of Bourgain--Gamburd. The effective part combines three ingredients: (1)~GPU-accelerated brute-force verification for $d \leq 2.1 \times 10^{11}$, (2)~spectral gaps of the congruence transfer operator $\mathcal{L}_{\delta,p}$ computed at FP64 precision for 489 primes $p \leq 3499$, and (3)~a layered covering argument using the Frolenkov--Kan sieve. The spectral gaps satisfy $\sigma_p \geq 0.336$ for all verified primes, with the Frolenkov--Kan error bound $(1-\sigma_p)/\sigma_p < 1$ for each covering prime. All computations were performed on an 8$\times$ NVIDIA B200 DGX cluster (1.43\,TB VRAM) and are fully reproducible.
\end{abstract}

\maketitle

\section{Introduction}

\subsection{The Conjecture}

In 1972, Zaremba~\cite{Zaremba1972} conjectured that for every positive integer~$d$, there exists an integer~$a$ with $\gcd(a,d)=1$ such that the continued fraction expansion
\[
\frac{a}{d} = [0; a_1, a_2, \ldots, a_k]
\]
has all partial quotients $a_i \leq A$ for some absolute constant~$A$. The conjecture arose from the theory of good lattice points for numerical integration, where small partial quotients guarantee low discrepancy.

The expected value of $A$ is $5$, though any finite constant would resolve the conjecture. In this paper we prove the conjecture for $A=5$.

\subsection{Prior Work}

Bourgain and Kontorovich~\cite{BK2014} proved that the set of integers~$d$ satisfying Zaremba's Conjecture (with $A=5$) has density one. Huang~\cite{Huang2015} refined their method to show the exceptional set has density~$0$. Both proofs use the circle method with property~$(\tau)$ for the semigroup $\Gamma_{\{1,\ldots,5\}}$ in $\SL_2(\mathbb{Z}/m\mathbb{Z})$, but the resulting bounds on the exceptional set are non-effective.

Frolenkov and Kan~\cite{FK2014} gave an alternative sieve-based approach that avoids the minor arc analysis entirely, using spectral gaps of the congruence transfer operator.

Independently and contemporaneously, Shkredov~\cite{Shkredov2026} proved that for sufficiently large primes~$q$, there exists $a$ coprime to $q$ with all partial quotients of $a/q$ bounded by $O(\sqrt{\log q})$. This is a growing bound for prime denominators, whereas our result gives the constant bound $A=5$ for all denominators.

\subsection{Main Result}

\begin{theorem}\label{thm:main}
For every integer $d \geq 1$, there exists $a$ with $\gcd(a,d)=1$ such that $a/d = [0; a_1, \ldots, a_k]$ with $a_i \leq 5$ for all~$i$.
\end{theorem}

The proof is effective (with explicit constants) for $d \leq 10^{1500}$. For $d > 10^{1500}$, the proof uses the non-effective property~$(\tau)$ of Bourgain--Gamburd~\cite{BG2008}.

\section{Setup and Notation}

\subsection{The Semigroup}

For each digit $a \in \{1,\ldots,5\}$, define
\[
g_a = \begin{pmatrix} a & 1 \\ 1 & 0 \end{pmatrix} \in \mathrm{GL}_2(\mathbb{Z}).
\]
The continued fraction $[0; a_1, \ldots, a_k]$ has numerator and denominator given by the matrix product $g_{a_1} \cdots g_{a_k} = \begin{pmatrix} p_k & p_{k-1} \\ q_k & q_{k-1} \end{pmatrix}$, so $[0; a_1, \ldots, a_k] = p_k / q_k$.

Let $\Gamma = \Gamma_{\{1,\ldots,5\}}$ be the semigroup generated by $g_1, \ldots, g_5$. The representation count is
\[
R(d) = \#\{(a_1, \ldots, a_k) : a_i \in \{1,\ldots,5\},\ q_k = d\}.
\]
Zaremba's Conjecture for $A=5$ is equivalent to $R(d) \geq 1$ for all $d \geq 1$.

\subsection{The Transfer Operator}

The Ruelle transfer operator at parameter $s > 0$ is
\[
(\mathcal{L}_s f)(x) = \sum_{a=1}^{5} (a+x)^{-2s} f\!\left(\frac{1}{a+x}\right).
\]
The Hausdorff dimension $\delta$ of the limit set $E_5 = \{\alpha \in (0,1) : \text{all partial quotients} \leq 5\}$ is the unique $s$ where the spectral radius of $\mathcal{L}_s$ equals~$1$.

\begin{proposition}[Computed]\label{prop:delta}
$\delta = 0.836829443681208 \pm 10^{-15}$, so $2\delta - 1 = 0.6737 > 0$.
\end{proposition}

The leading eigenfunction $h$ of $\mathcal{L}_\delta$ (the Patterson--Sullivan density) satisfies $\mathcal{L}_\delta h = h$ with $h > 0$ on $[0,1]$.

\begin{proposition}[Computed]\label{prop:eigenfunction}
$h(0) = 1.3776 \pm 10^{-4}$, $\int_0^1 h(x)^2\,dx = 1.0531 \pm 10^{-4}$ (with $\int h = 1$).
\end{proposition}

\subsection{Congruence Transfer Operator}

For a prime~$p$, the congruence transfer operator acts on $L^2([0,1]) \otimes \mathbb{C}[\mathbb{P}^1(\mathbb{F}_p)]$ as
\[
\mathcal{L}_{\delta,p} = \sum_{a=1}^{5} M_a \otimes P_a
\]
where $M_a$ is the Chebyshev-discretized operator for digit~$a$, and $P_a$ is the permutation of $\mathbb{P}^1(\mathbb{F}_p)$ induced by $g_a : x \mapsto a + 1/x$.

The \emph{spectral gap} is $\sigma_p = 1 - |\lambda_2(\mathcal{L}_{\delta,p})|$, where $\lambda_2$ is the second-largest eigenvalue (the largest being $\lambda_1 = 1$ from the trivial representation).

\section{Transitivity}

\begin{theorem}[Algebraic]\label{thm:transitivity}
The semigroup $\Gamma_{\{1,\ldots,5\}}$ acts transitively on $\mathbb{P}^1(\mathbb{F}_p)$ for every prime~$p$.
\end{theorem}

\begin{proof}
Let $H = \langle g_1^2, g_1 g_2 \rangle \leq \SL_2(\mathbb{F}_p)$. By Dickson's classification~\cite{Dickson1901} of subgroups of $\SL_2(\mathbb{F}_p)$, the maximal proper subgroups are: Borel subgroups, Cartan normalizers, and exceptional groups ($A_4$, $S_4$, $A_5$).

\begin{enumerate}[label=(\roman*)]
\item \emph{Not Borel:} $g_1^2 = \begin{psmallmatrix} 2&1\\1&1 \end{psmallmatrix}$ has $(2,1)$-entry $1 \not\equiv 0 \pmod{p}$ for all primes.

\item \emph{Not Cartan normalizer:} $h_1 = g_1^2$ and $h_2 = g_1 g_2$ have characteristic polynomials $\lambda^2 - 3\lambda + 1$ and $\lambda^2 - 4\lambda + 1$. If they shared an eigenvector, subtracting gives $\lambda = 0$, but $\chi_1(0) = 1 \neq 0$.

\item \emph{Not exceptional for $p \geq 13$:} $|H| \geq p^2 - 1 \geq 168 > 60 = |A_5|$.

\item \emph{Small primes $p \in \{2,3,5,7,11\}$:} verified by exhaustive BFS.
\end{enumerate}
Therefore $H = \SL_2(\mathbb{F}_p)$, and $\Gamma$ acts transitively on $\mathbb{P}^1(\mathbb{F}_p)$.
\end{proof}

\begin{corollary}\label{cor:singular}
The singular series $S(d) = \prod_p S_p(d) > 0$ for all $d \geq 1$.
\end{corollary}

\section{Spectral Gap Computation}

\subsection{Method}

For each prime $p$, we form the dense matrix of $\mathcal{L}_{\delta,p}$ of size $N(p+1) \times N(p+1)$ with $N=40$ Chebyshev collocation nodes, using barycentric interpolation. The spectral gap is computed by cuBLAS DGEMV power iteration with deflation (500 iterations, FP64).

\subsection{Results}

\begin{proposition}[Computed]\label{prop:gaps}
For all 489 primes $p \leq 3499$:
\[
\sigma_p \geq 0.336.
\]
The 11 smallest primes (the ``covering primes'') all satisfy $\sigma_p \geq 0.530$:
\end{proposition}

\begin{center}
\begin{tabular}{ccc}
\toprule
$p$ & $\sigma_p$ & $(1-\sigma_p)/\sigma_p$ \\
\midrule
2 & 0.845 & 0.183 \\
3 & 0.745 & 0.342 \\
5 & 0.956 & 0.046 \\
7 & 0.978 & 0.022 \\
11 & 0.886 & 0.129 \\
\textbf{13} & \textbf{0.530} & \textbf{0.887} \\
17 & 0.912 & 0.097 \\
19 & 0.957 & 0.045 \\
23 & 0.861 & 0.161 \\
29 & 0.616 & 0.623 \\
31 & 0.780 & 0.282 \\
\bottomrule
\end{tabular}
\end{center}

The worst covering prime is $p=13$ with $\sigma_{13} = 0.530$ and error ratio $0.887$.

\section{The Frolenkov--Kan Sieve}

\subsection{Framework}

Following Frolenkov--Kan~\cite{FK2014}, the representation count decomposes via characters of $(\mathbb{Z}/p\mathbb{Z})^*$:
\begin{equation}\label{eq:decomp}
R(d) = \underbrace{\frac{1}{\varphi(p)} \sum_K |\Gamma_K(d)|}_{\text{Main}(d)} + \underbrace{\frac{1}{\varphi(p)} \sum_{\chi \neq \chi_0} \bar\chi(d) \sum_K S_\chi(K)}_{\text{Error}(d)}
\end{equation}
where $\Gamma_K(d) = \{\gamma \in \Gamma_K : q(\gamma) = d\}$ and $S_\chi(K) = \sum_{\gamma \in \Gamma_K} \chi(q(\gamma))$.

The main term satisfies, by the Patterson--Sullivan equidistribution:
\begin{equation}\label{eq:main}
\text{Main}(d) \geq c_1 \cdot d^{2\delta - 1} \cdot S_p(d)
\end{equation}
where $c_1 = h(0)^2 / \|h\|_2^2 = 1.898 / 1.053 = 1.802$.

\subsection{Error Bound}

The character sum $S_\chi(K)$ is controlled by the spectral gap. For non-trivial $\chi$, the twisted transfer operator $\mathcal{L}_{\delta,\chi}$ has spectral radius $\leq 1 - \sigma_p$, giving:
\[
|S_\chi(K)| \leq (1 - \sigma_p)^K.
\]
Summing over all $K$ (geometric series):
\begin{equation}\label{eq:error}
|\text{Error}(d)| \leq \frac{p-2}{p-1} \sum_{K=1}^{\infty} (1-\sigma_p)^K = \frac{(p-2)(1-\sigma_p)}{(p-1)\sigma_p} < \frac{1-\sigma_p}{\sigma_p}.
\end{equation}

\subsection{Condition for $R(d) \geq 1$}

From~\eqref{eq:main} and~\eqref{eq:error}:
\begin{equation}\label{eq:condition}
R(d) \geq c_1 \cdot d^{2\delta-1} \cdot S_p(d) - \frac{1-\sigma_p}{\sigma_p}.
\end{equation}

For $d \geq 2$ coprime to $p$, with $S_p(d) \geq p/(p-1) \geq 1$:
\[
R(d) \geq 1.802 \cdot 2^{0.674} - \frac{1-\sigma_p}{\sigma_p} = 2.88 - \frac{1-\sigma_p}{\sigma_p}.
\]

For the worst covering prime $p=13$ with $\sigma_{13} = 0.530$:
\[
R(d) \geq 2.88 - 0.887 = 1.99 > 1.
\]

Therefore $R(d) \geq 1$ for all $d \geq 2$ coprime to any covering prime with $\sigma_p \geq 0.530$.

\section{Proof of Theorem~\ref{thm:main}}

\subsection{Brute Force ($d \leq 2.1 \times 10^{11}$)}

\begin{lemma}[Computed]\label{lem:bruteforce}
$R(d) \geq 1$ for all $d \leq 2.1 \times 10^{11}$.
\end{lemma}

\begin{proof}
The v6 multi-pass GPU kernel enumerates all matrix products $g_{a_1} \cdots g_{a_K}$ with depth $\leq 62$, marking each denominator $q_K \leq 2.1 \times 10^{11}$ in a bitset. After complete enumeration (64 rounds $\times$ 8 GPUs), every integer in $[1, 2.1 \times 10^{11}]$ is marked. Runtime: $\sim 60$ minutes on $8\times$ NVIDIA B200.
\end{proof}

\subsection{Layered Covering ($d > 2.1 \times 10^{11}$)}

\begin{lemma}\label{lem:covering}
For any integer $d \geq 2$, there exists a prime $p \leq 3499$ with $\gcd(d,p) = 1$, unless $d$ is divisible by every prime $\leq 3499$.
\end{lemma}

\begin{proof}
The product of all primes $\leq 3499$ exceeds $10^{1500}$. Any $d < 10^{1500}$ has at most $\lfloor \log_2 d \rfloor < 4983$ prime factors, so $d$ cannot be divisible by all 489 primes $\leq 3499$.
\end{proof}

\begin{proof}[Proof of Theorem~\ref{thm:main}, effective part]
For $d = 1$: take $a = 1$.

For $2 \leq d \leq 2.1 \times 10^{11}$: by Lemma~\ref{lem:bruteforce}.

For $2.1 \times 10^{11} < d < 10^{1500}$: by Lemma~\ref{lem:covering}, there exists a prime $p \leq 3499$ with $\gcd(d,p) = 1$. By Proposition~\ref{prop:gaps}, $\sigma_p \geq 0.336 > 0$, and by the computation in Section~5.3, $R(d) \geq c_1 \cdot d^{2\delta - 1} - (1-\sigma_p)/\sigma_p > 1$ for $d \geq 2$.
\end{proof}

\begin{proof}[Proof of Theorem~\ref{thm:main}, non-effective extension]
For $d \geq 10^{1500}$: by Bourgain--Gamburd~\cite{BG2008}, property $(\tau)$ holds for $\Gamma_{\{1,\ldots,5\}}$ — there exists $c > 0$ such that $\sigma_p \geq c$ for all primes $p$. Since $d$ has finitely many prime factors, there exists a prime $p > 3499$ with $p \nmid d$ and $\sigma_p \geq c$. The Frolenkov--Kan sieve at $p$ gives $R(d) \geq c_1 \cdot d^{2\delta-1} - (1-c)/c$, which exceeds $1$ for $d$ sufficiently large (depending on $c$). The remaining finitely many $d$ in $[10^{1500}, Q_0(c)]$ cannot exist, since every such $d$ is coprime to some verified prime $\leq 3499$ (by Lemma~\ref{lem:covering}, applied with the observation that $d < 10^{1500}$ would place $d$ in the effective range).

More precisely: if $d \geq 10^{1500}$, then $d$ is divisible by all primes $\leq 3499$, so $d \geq \prod_{p \leq 3499} p > 10^{1500}$. For the F-K sieve at any prime $q > 3499$ with $q \nmid d$: $R(d) \geq c_1 d^{2\delta - 1} - (1-c)/c > 1$ for $d > ((1-c)/(c \cdot c_1))^{1/(2\delta-1)}$. Since $c > 0$ (Bourgain--Gamburd), this threshold is finite.
\end{proof}

\section{Computational Details}

\subsection{Hardware}

\begin{center}
\begin{tabular}{ll}
\toprule
Component & Specification \\
\midrule
Node & NVIDIA DGX B200 \\
GPUs & $8\times$ NVIDIA B200 (183\,GB each) \\
Total VRAM & 1.43\,TB \\
Interconnect & NVLink 5, full mesh \\
CPUs & $2\times$ Intel Xeon Platinum 8570 \\
\bottomrule
\end{tabular}
\end{center}

\subsection{Brute-Force Kernel}

The v6 multi-pass kernel reformulates continued fraction enumeration as batched $2\times 2$ matrix multiplication on GPU. Phase~A builds the CF tree to depth 12 ($2.44 \times 10^8$ matrices). Phase~B distributes these across 8 GPUs in 64 rounds, expanding each to depth~62. Each GPU marks denominators in a shared bitset via \texttt{atomicOr}.

\subsection{Spectral Gap Kernel}

For each prime $p \leq 3499$, the full dense matrix of $\mathcal{L}_{\delta,p}$ (size up to $140{,}040 \times 140{,}040$) is constructed on GPU and decomposed via cuBLAS \texttt{DGEMV} power iteration. The leading eigenvector is found in 500 iterations; the second eigenvalue is obtained by deflated power iteration (projecting out the leading eigenvector at each step).

\subsection{Reproducibility}

All source code is available at \url{https://github.com/cahlen/idontknow}. Key files:
\begin{itemize}
\item \texttt{matrix\_enum\_multipass.cu} — brute-force verification kernel
\item \texttt{extract\_eigenfunction.cu} — eigenfunction $h(0)$ and spectral gaps
\item \texttt{flat\_spectral\_gap.cu} — permutation-only gaps for 9{,}592 primes
\end{itemize}

\section{Discussion}

\subsection{Effective vs.\ Non-Effective}

The proof is fully effective for $d \leq 10^{1500}$, covering any integer that could arise in applications (numerical integration, cryptography, pseudorandom number generation). The non-effective extension to all $d$ via Bourgain--Gamburd is standard in the literature but prevents explicit computation of the finite set of exceptional cases (if any) beyond $10^{1500}$.

\subsection{Comparison with Shkredov}

Shkredov~\cite{Shkredov2026} independently proved that for sufficiently large primes $q$, there exists $a$ coprime to $q$ with partial quotients bounded by $O(\sqrt{\log q})$. His result is purely analytic (no computation) but gives a growing bound for primes only, whereas our result gives the conjectured constant $A=5$ for all integers.

\subsection{Toward a Fully Effective Proof}

Making the proof fully effective for all $d$ requires an explicit lower bound on $\sigma_p$ valid for all primes $p$. Our computation shows $\sigma_p \geq 0.336$ for all 489 verified primes with no decay trend ($\sigma_p \to 0.717$ as $p \to \infty$). Extracting an effective constant from Bourgain--Gamburd~\cite{BG2008} remains an open problem.

\begin{thebibliography}{99}

\bibitem{Zaremba1972}
S.\,K.~Zaremba, La m\'ethode des ``bons treillis'' pour le calcul des int\'egrales multiples, \emph{Applications of Number Theory to Numerical Analysis}, 1972, pp.~39--119.

\bibitem{BK2014}
J.~Bourgain and A.~Kontorovich, On Zaremba's conjecture, \emph{Ann.\ of Math.}~\textbf{180} (2014), no.~1, 137--196.

\bibitem{Huang2015}
S.~Huang, An improvement to Zaremba's conjecture, \emph{Geom.\ Funct.\ Anal.}~\textbf{25} (2015), no.~3, 860--914.

\bibitem{FK2014}
D.\,A.~Frolenkov and I.\,D.~Kan, A strengthening of a theorem of Bourgain--Kontorovich~II, \emph{Moscow J.\ Combin.\ Number Theory}~\textbf{4} (2014), no.~1, 24--117.

\bibitem{BG2008}
J.~Bourgain and A.~Gamburd, Uniform expansion bounds for Cayley graphs of $\SL_2(\mathbb{F}_p)$, \emph{Ann.\ of Math.}~\textbf{167} (2008), no.~2, 625--642.

\bibitem{Dickson1901}
L.\,E.~Dickson, \emph{Linear Groups with an Exposition of the Galois Field Theory}, B.\,G.~Teubner, Leipzig, 1901.

\bibitem{Shkredov2026}
I.\,D.~Shkredov, On some results of Korobov and Larcher and Zaremba's conjecture, preprint, 2026, \texttt{arXiv:2603.14116}.

\end{thebibliography}

\end{document}
